\chapter{Local fields, sewing overlaps, and Ward tests}
\label{ch:overlaps}

Energy levels tell us which states might exist; matrix elements tell us how
their fields interact. This chapter develops two ways to approach those
interactions from a lattice chain. One uses fixed local operators and their
matrix elements. The other glues normalized wavefunctions into a
pair-of-pants geometry. Both are natural consequences of the state-operator
correspondence, but they have different finite-size corrections. We derive
phase-independent three-point invariants, the unavoidable descendant content
of a spatially averaged primary, and a useful response test for a proposed
running coupling. These calculations also explain why a large overlap with
one low state is not a general error bound for an OPE coefficient.
The primary and descendant conventions are those of
Chapter~\ref{ch:cft}; Chapter~\ref{ch:sewing} explains the full-CFT
consistency conditions that any limiting overlap interpretation must obey.

\section{From a local lattice source to a normalized field}
\index{operator calibration}

Let $o_j$ be a fixed finite-range lattice density. Its momentum-filtered
average is
\[
 O_k=\frac1L\sum_j e^{-ikj}o_j.
\]
For a neutral source its vacuum mean should first be subtracted. For a
field at a nonzero carrier momentum, the phase removes that carrier before
comparison with a residual scalar zero mode. The support and coefficients
of $o_j$ should be kept fixed while $L$ changes.

Suppose that in a critical limit
\[
 o_j\sim a_\phi\phi(x)+\sum_ba_b\mathcal O_b(x).
\]
The coefficient $a_\phi$ is a nonuniversal microscopic normalization.
If $|\phi\rangle$ is the corresponding unit-normalized cylinder state,
then
\[
 A_\phi(L)=\langle\phi|O|0\rangle
       \sim a_\phi(2\pi/L)^{\Delta_\phi}
\]
in a fixed coordinate convention. A ratio such as
\begin{equation}
 c_O(L)=\frac{\langle\phi^\dagger|O|\phi\rangle}
                    {\langle\phi|O|0\rangle}
 \label{eq:lattice-cubic}
\end{equation}
can cancel $a_\phi$ and the common cylinder factor. Under convergence to
the intended primary insertion and states, it approaches the normalized
three-point coefficient $C(\phi,\phi,\phi)$.

The conditional phrase matters. The denominator tests $O$ on the vacuum;
the numerator tests it on an excited state. An unwanted operator can have
a small vacuum-to-low-state matrix element and a substantial excited-state
matrix element. No algebraic inequality turns a single large vacuum
overlap into a uniform operator-norm estimate.

\subsection{A finite-dimensional operator optimization}

One can choose a span of local words $o^{(a)}$ with the desired charge and
momentum and set $O=\sum_ac_aO^{(a)}$. The source Gram matrix is
\[
 G_{ab}=\langle0|(O^{(a)})^\dagger O^{(b)}|0\rangle.
\]
The source squared norm is $c^\dagger Gc$. Linear constraints can normalize the
overlap with a target and cancel overlaps with selected unwanted states.
Minimizing the Gram norm subject to those constraints is a finite linear
algebra problem. Null operator directions must be removed consistently;
otherwise a formally large coefficient can represent a zero or nearly
zero source.

The resulting recipe is useful when trained at one reference size and
then transferred unchanged. Retraining independently at every size may
produce a sequence of operators with no fixed local continuum limit.
Equally, choosing the recipe that most closely matches a favorite CFT
value makes subsequent agreement circular. Predetermined operator ranges,
all successful recipes, and failures should remain visible.

\begin{example}[Solving the source optimization]
Assume the independent source Gram matrix is positive definite,
and define $b_a=\langle0|(O^{(a)})^\dagger|\phi\rangle$.
The target amplitude of $O=\sum_ac_aO^{(a)}$ is $b^\dagger c$.
For $b\ne0$, minimizing $c^\dagger Gc$ subject to
$b^\dagger c=1$ gives
\[
 c_{\mathrm{opt}}=\frac{G^{-1}b}{b^\dagger G^{-1}b},\qquad
 \|O_{\mathrm{opt}}|0\rangle\|^2
 =\frac1{b^\dagger G^{-1}b}.
\]
Indeed Cauchy--Schwarz applied to $G^{1/2}c$ and $G^{-1/2}b$
gives the lower bound, saturated by the displayed vector. The
maximum target fraction is $b^\dagger G^{-1}b$, the squared
projection of $|\phi\rangle$ onto the allowed source span.

For two normalized sources with real mutual overlap $1/5$ and
target overlaps $3/5$ each, $G$ has diagonal entries one and
off-diagonal entries $1/5$, while $b=(3/5,3/5)^T$.
The optimum is $c=(5/6,5/6)^T$ and its fraction is $3/5$.
Each source alone has fraction $9/25$. The improvement is a
finite linear-algebra fact about the chosen span, with no CFT
assumption. Whether this optimized source has a stable local
continuum interpretation is the separate transfer-in-size test.
\end{example}

\begin{definition}[Vacuum-source fraction]
For a normalized target $|\phi\rangle$ and source $s=O|0\rangle$, define
\[
 f_\phi=\frac{|\langle\phi|s\rangle|^2}{\|s\|^2}.
\]
The complement $1-f_\phi$ is exactly the fraction of the finite source
squared norm outside that state, provided the norm and normalization are known.
\end{definition}

This is a strong finite statement, even if no other states have been
saved. What it does not measure is $\|O|a\rangle\|$ for other $|a\rangle$
or the difference between $O$ and a continuum operator in a norm valid on
an entire low-energy space.

\section{Complex phases and information that cannot be rephased away}
\index{phase invariant!OPE coefficients}

Normalize charged fields by
$\langle\phi(z)\phi^\dagger(0)\rangle=|z|^{-2\Delta_\phi}$.
The transformation $\phi\mapsto e^{i\theta_\phi}\phi$ preserves this
normalization and changes the cubic $c_\phi=C(\phi,\phi,\phi)$ by
$e^{3i\theta_\phi}$. Thus an individual complex cubic has a conventional
phase.

For the charge pattern
\[
 \phi:(1,0),\quad\chi:(0,1),\quad
 \psi:(1,1),\quad\eta:(1,-1)
\]
under $C_3\times C_3$, define
\[
 \lambda=C(\phi,\chi,\psi^\dagger),\quad
 \nu=C(\phi,\psi,\eta),\quad
 \kappa=C(\chi,\psi,\eta^\dagger).
\]
The first charge is intrinsic alpha charge. The second records the
proposed order-three internal part of the carrier translation,
distinct from the residual conformal spin. Its survival as a faithful
continuum symmetry is a premise tested in Chapter~\ref{ch:identification}.
The rephasing algebra below uses the displayed selection rules; it
does not by itself prove that continuum symmetry interpretation.
All three are allowed by charge conservation. Their phases transform as
\[
 \lambda\mapsto e^{i(\theta_\phi+\theta_\chi-\theta_\psi)}\lambda,
 \quad
 \nu\mapsto e^{i(\theta_\phi+\theta_\psi+\theta_\eta)}\nu,
\]
and similarly for $\kappa$. Cubing and dividing by self cubics cancels
the field phases:
\begin{equation}
 I_\lambda=\frac{\lambda^3c_\psi}{c_\phi c_\chi},\qquad
 I_\nu=\frac{\nu^3}{c_\phi c_\psi c_\eta},\qquad
 I_\kappa=\frac{\kappa^3c_\eta}{c_\chi c_\psi}.
 \label{eq:triangle-invariants}
\end{equation}
The ratio $J=\kappa\nu/(\lambda c_\psi)$ is also invariant and obeys
$J^3=I_\nu I_\kappa/I_\lambda$.

\begin{example}[Parity need not make every invariant real]
Suppose spatial parity acts as
$P\phi=\phi$, $P\chi=\chi^\dagger$,
$P\psi=\eta$, $P\eta=\psi$, in a suitable phase convention.
Then $c_\eta=c_\psi$ and parity can make $\kappa$ and $c_\chi$ real,
so $I_\kappa$ is real. But parity sends $\lambda$ to the different
coupling $C(\phi,\chi^\dagger,\eta^\dagger)$. It does not send every
field to its adjoint. It therefore does not force $I_\lambda$ to equal
its complex conjugate.
\end{example}

A unitary symmetry that simultaneously conjugates all three fields in
$I_\lambda$ would force its reality: its field phases cancel between
numerator and denominator. A nonreal invariant can therefore obstruct
such a symmetry. This is more informative than finding a complex entry
in an arbitrary eigenvector basis. It is also different from the
antiunitary conjugation associated with a real Hamiltonian.

\begin{numericalrecord}[Finite mixed-tensor fingerprint]
At $L=15$, prescribed range-four/five recipes give
$|\lambda|\simeq1.420$--$1.432$. Three different choices of which
field supplies the inserted local operator agree within approximately
$.26$--$.68\%$ for matched recipes. The direct chi insertion gives
\[
 \Re I_\lambda\in[1.114433,1.119331],\qquad
 \Im I_\lambda\in[-.203273,-.199851].
\]
These ranges describe deterministic operator variation. They are not
confidence intervals for continuum data. The nonreal finite invariant
is nevertheless unchanged by the saved state rephasings.
\end{numericalrecord}
Chapter~\ref{ch:identification} applies this invariant to a specified
canonical Cardy hypothesis, with the required spectral completeness
and sewing assumptions kept explicit.

\section{A local primary creates an entire translation family}
\index{descendant!global Virasoro}

The state-operator correspondence gives a precise benchmark for a source
fraction. Let $\phi$ be a scalar primary with $h=\bar h=\Delta/2$.
On the plane,
\[
 \phi(z,\bar z)|0\rangle
 =e^{zL_{-1}+\bar z\bar L_{-1}}|\phi\rangle.
\]
Radial quantization and the global conformal algebra imply
\[
 \|L_{-1}^n|h\rangle\|^2=n!(2h)_n,
\]
where $(a)_n=a(a+1)\cdots(a+n-1)$. This follows inductively from
$[L_1,L_{-1}]=2L_0$.

On the unit circle, spatial averaging keeps precisely the terms with
equal left and right levels. The normalized scalar translation
descendant is
\[
 |D_n\rangle=
 \frac{L_{-1}^n\bar L_{-1}^n|\phi\rangle}{n!(\Delta)_n}.
\]
Its amplitude relative to the primary amplitude is $(\Delta)_n/n!$.
Orthogonality of different energies then gives
\begin{equation}
 F_{\mathrm{vac}}(\Delta)=
 \frac{\|O_0|0\rangle\|^2}{|\langle\phi|O_0|0\rangle|^2}
 =\sum_{n=0}^{\infty}\left[\frac{(\Delta)_n}{n!}\right]^2.
 \label{eq:vac-descendant-series}
\end{equation}
For $0\leq\Delta<1/2$, Gauss summation yields
\begin{equation}
 F_{\mathrm{vac}}(\Delta)
 =\frac{\Gamma(1-2\Delta)}{\Gamma(1-\Delta)^2},\qquad
 p_{\mathrm{primary}}=
 \frac{\Gamma(1-\Delta)^2}{\Gamma(1-2\Delta)}.
 \label{eq:primary-Gamma-fraction}
\end{equation}
The first daughter has relative amplitude $\Delta$ and squared weight
$\Delta^2$. A pure local primary therefore creates more than its
primary state. Demanding a source fraction of exactly one would impose
the wrong continuum target.

The convergence condition has a physical meaning. The equal-time
two-point function behaves near coincidence as $|\theta|^{-2\Delta}$.
Its spatial integral is finite precisely for $\Delta<1/2$. For larger
dimensions, a microscopic contact region can dominate the unfiltered
source norm even though its low-energy matrix elements have a good
scaling limit.

\subsection{Why a finite fraction can exceed the CFT fraction}

Let $O=L^{-1}\sum_jo_j$. A regular short-distance contribution to its
variance is of order $L^{-1}$, whereas the long-distance primary
contribution is of order $L^{-2\Delta}$. Thus a possible expansion is
\[
 \operatorname{Var}(O)=A L^{-2\Delta}+B/L+\cdots.
\]
After normalization by the primary overlap squared,
\begin{equation}
 F_L=F_{\mathrm{vac}}+B'L^{2\Delta-1}+\cdots.
 \label{eq:contact-correction}
\end{equation}
The coefficient $B'$ is not constrained to be positive: it measures
the difference between a microscopic regularization and the ideal
continuum contribution. At light dimensions near $.13$--$.28$, this
correction decays much more slowly than $L^{-2}$.

For example, using the conditional $L=18$ stress dimensions gives
primary fractions approximately $.95865,.79759,.84008$ for
$\phi,\epsilon,\chi$. Measured fixed range-five source fractions are
approximately $.985,.962,.969$. These numbers do not prove that the
continuum fields are impure or that genuine descendants have been
removed. They show why finite source purity cannot be turned into a
percent-level continuum OPE error merely by subtracting from one.

\subsection{Excited sources have a different descendant coefficient}

There is a related but different calculation for $O_0|\phi\rangle$.
In this calculation $O_0$ is the unit-circle average of the normalized
continuum primary. A microscopic source must first have its common
vacuum-to-primary amplitude divided out before making this comparison.
Suppose a scalar primary $\beta$ of weights $(h,h)$ occurs in
$\phi\times\phi$, with coefficient $C_\beta$. For equal external
weights, the chiral global OPE coefficient multiplying
$L_{-1}^n\beta$ is $(h)_n/[n!(2h)_n]$. After normalizing left and right
descendants, the scalar amplitude is
\[
 C_\beta\frac{(h)_n^2}{n!(2h)_n}.
\]
Consequently
\begin{equation}
 \|O_0|\phi\rangle\|^2\geq
 \sum_{\beta\ \mathrm{scalar}}|C_\beta|^2
 \sum_{n\geq0}\frac{(h_\beta)_n^4}
 {(n!)^2(2h_\beta)_n^2}.
 \label{eq:excited-descendant-bound}
\end{equation}
Other Virasoro directions and higher-spin primary families can add
positive weight. The first scalar daughter amplitude here is
$C_\beta\Delta_\beta/4$, whereas the vacuum-source amplitude above
was $\Delta$. Mixing these two external-state geometries produces a
factor-of-four error.

If the identity is allowed in the channel, its global family contributes
only the $n=0$ term. One should not substitute $h=0$ into the displayed
ratio of Pochhammer symbols for $n>0$: the vacuum's translation
descendants vanish. Other vacuum Virasoro directions, such as
$T\bar T$, are separate positive contributions.

\section{Ward vectors, phases, and nonorthogonal candidate spans}

A candidate daughter should satisfy energy, momentum, norm, and matrix
element tests together. A lattice approximation to a Virasoro mode can
produce a Ward vector $S$. One may project it onto a fixed collection
of residual-checked low states, but those states need not be exactly
orthogonal after numerical reconstruction.

Let the columns of $Y$ be the retained states, let
$G=Y^\dagger Y$, and put $B=Y^\dagger S$. For positive invertible
$G$, the actual orthogonal projector is
\[
 \Pi_Y=Y G^{-1}Y^\dagger,
 \qquad \|\Pi_YS\|^2=B^\dagger G^{-1}B.
\]
If $A=Y^\dagger O|0\rangle$, the overlap with the coherently phased
normalized Ward projection is
\[
 \frac{B^\dagger G^{-1}A}
 {\sqrt{B^\dagger G^{-1}B}\,\langle\phi|O|0\rangle}.
\]
Independent rephasings of the numerical eigenvectors cancel in this
formula. Replacing the complex amplitudes by their absolute values
before forming the sum destroys that coherent test.

In an earlier calculation, a two-candidate span at $L=15$
has about $.9870$ of a conditional ideal daughter squared norm, yet only
$.2418$ of the complete Ward-source squared norm. Its bare-source vacuum
ratios have real parts near $.098$, compared with the conditional
primary dimension $.1464$; trained range-five ratios are smaller.
The attractive norm agreement therefore does not settle the family
assignment. An energy near ``primary plus two'' is only one of several
necessary checks.

\begin{example}[A local improvement preserves an equal-energy numerator]
For a fixed local operator $O$ and real $t$, define
\[
 O_t=O+t[H,[H,O]].
\]
Between exact energy eigenstates,
\[
 \langle n|O_t|m\rangle
 =\bigl[1+t(E_n-E_m)^2\bigr]\langle n|O|m\rangle.
\]
If $\phi$ and $\phi^\dagger$ are exactly degenerate, the cubic
numerator in \eqref{eq:lattice-cubic} is unchanged. With primary and
daughter gaps $g_0,g_1$, their vacuum ratio changes by
$(1+tg_1^2)/(1+tg_0^2)$. Choosing $t=-g_1(L_\star)^{-2}$ kills the
daughter amplitude at one reference size while changing the primary
denominator much less when $g_0\ll g_1$.

The same fixed $t$ must then be used at every size. If the gaps scale
as $L^{-1}$, the improvement vanishes relatively as $L^{-2}$ in the
limit. Retuning $t$ with $L$ would be a different operator sequence.
This example does not assert that the actual diagonal word recipes
have this commutator form; it exhibits a concrete possible mechanism.
\end{example}

\section{Heat filtering and a conditional omitted-tail bound}
\index{heat filtering}

A positive-time filter reduces sensitivity to short-distance source
weight. For exact finite eigenstates define
\[
 F_L(\tau)=
 \frac{\|e^{-\tau(H-E_\phi)}O|0\rangle\|^2}{|A_\phi|^2}
 =\sum_j\left|\frac{A_j}{A_\phi}\right|^2
       e^{-2\tau(E_j-E_\phi)}.
\]
The candidate primary contributes one. In an ideal CFT, with
$\omega_L=2\pi v/L$, the scalar descendants lie $2n\omega_L$ above
it, giving
\begin{equation}
 F_{\mathrm{CFT}}(\tau)=
 {}_2F_1(\Delta,\Delta;1;q),\qquad
 q=e^{-4\tau\omega_L}.
 \label{eq:filtered-hypergeometric}
\end{equation}
Keeping $\tau\omega_L$ fixed as $L$ grows suppresses microscopic
energies while preserving nontrivial descendant weights. Keeping
physical $\tau$ fixed sends $q\to1$ instead.

Suppose a set $J$ of exact orthonormal eigenstates is retained and every
omitted source-coupled state has energy at least $E_{\rm cut}$. Positivity
then proves
\[
 0\leq F_L(\tau)-F_J(\tau)
 \leq e^{-2\tau(E_{\rm cut}-E_\phi)}
 \left[F_L(0)-\sum_{j\in J}|A_j/A_\phi|^2\right].
\]
The full unfiltered norm is essential. So is the true omitted-energy
cutoff: the largest energy among a few accurate Ritz vectors is not
automatically such a cutoff. This distinction separates a rigorous
positive bound from a useful conditional spectral diagnostic.

\section{Sewing two circles into one}
\index{wavefunction overlap}

Let $L=L_1+L_2$. In a common height embedding, define the literal
wavefunction overlap
\[
 A_{abc}(L_1,L_2)=
 \langle c,L|\bigl(|a,L_1\rangle\otimes|b,L_2\rangle\bigr).
\]
The outgoing parent is a bra. Seams, orientation, fusion-vertex gauge,
and normalization of all three states must be held fixed. Exact
microscopic gluing is a well-defined finite operation even when its
continuum description has not been established.

\subsection{What the finite joining map actually does}
\label{sec:overlap-literal-joining}

The tensor product in the overlap formula is taken in the ordinary
height embedding of Chapter~\ref{ch:chain}. It does not mean that two
closed fusion paths automatically form a closed path of the combined
length. Write
\[
 \iota_L:\HH_{\mathrm{path},L}\longrightarrow(\C^6)^{\otimes L}
\]
for the isometric embedding that sends each admissible path to its
height word. Its adjoint deletes the forbidden-word components.
After identifying the two ambient tensor products with the ambient
length-$L$ space, the actual joining map is
\begin{equation}
 J_{L_1,L_2}
 =\iota_L^\dagger(\iota_{L_1}\otimes\iota_{L_2}),
 \qquad L=L_1+L_2.
 \label{eq:literal-joining-map}
\end{equation}
Thus
\[
 A_{abc}
 =\langle c,L|J_{L_1,L_2}
             \bigl(|a,L_1\rangle\otimes|b,L_2\rangle\bigr).
\]
The children and parent are individually normalized before this
contraction.

To see the seam constraint, let $u$ and $v$ be allowed child words.
The concatenated word retains all their internal edges, but replaces
their two closing edges by
\[
 u_{L_1-1}\longrightarrow v_0,\qquad
 v_{L_2-1}\longrightarrow u_0.
\]
The image under $J$ is the concatenated parent basis vector if both
new edges are allowed, and zero otherwise. For example, the length-three
words $(I_0,X_0,X_0)$ and $(I_1,X_1,X_1)$ are each valid closed
paths. Their concatenation fails both new seam constraints:
$X_0\rho$ does not contain $I_1$, and $X_1\rho$ does not contain
$I_0$. Its image is zero. Concatenating $(I_0,X_0,X_0)$ with
itself instead gives a valid unit-norm parent word.

Distinct surviving child pairs give distinct parent words. Hence
$J^\dagger J$ is the orthogonal projector onto the child pairs that
survive the two seam tests; $JJ^\dagger$ projects onto the parent
words whose two cut blocks separately close. In particular $J$ is
a partial isometry: its nonzero singular values are one, and
\[
 \|J\|\leq1,\qquad
 |A_{abc}|\leq
 \bigl\|J(|a\rangle\otimes|b\rangle)\bigr\|\leq1
\]
for normalized states. One should not renormalize the projected joined
vector separately for each pair of children. That would change the
multilinear overlap tensor being compared with the pants amplitude.

This elementary contraction property is also the source of the
error bound used later. If $u,v$ and $\widetilde u,\widetilde v$
are normalized child states with respective norm differences
$\delta_1,\delta_2$, then
\[
 \|u\otimes v-\widetilde u\otimes\widetilde v\|
 \leq\delta_1+\delta_2.
\]
Contracting with $J$ and a fixed unit parent cannot increase this
bound. Dividing both amplitudes by the same nonzero vacuum overlap
therefore gives a difference at most
$(\delta_1+\delta_2)/|A_{000}|$. A small vacuum overlap can amplify
an otherwise small absolute state error.

\paragraph{Categorical cups use a different normalization.}
There is no extra quantum-dimension factor in
\eqref{eq:literal-joining-map}: its height basis is orthonormal.
Square roots of quantum dimensions do occur when a graphical
categorical contraction is converted to normalized states. For a
simple object $x$, a standard unitary coevaluation
$R_x:\one\to x^*\otimes x$ obeys
\[
 R_x^\dagger R_x=d_x\,1_{\one}.
\]
The cup vector of norm one is consequently $d_x^{-1/2}R_x$.
For $x=\rho$, this uses the positive dimension
$d=(3+\sqrt{13})/2$. The unnormalized closed cup--cap diagram
evaluates to $d_x$, while the norm of its normalized cup is one.
These statements concern categorical duality maps; they do not
insert a universal $\sqrt{d_x}$ into a literal height overlap.
A different graphical sewing must specify its duality maps and
their conversion to the chosen orthonormal basis.

The same distinction applies to gauge and phase choices. Under
unitary changes of path coordinates $U_L$, representing the same
finite map requires
\[
 J\longmapsto
 U_LJ(U_{L_1}^\dagger\otimes U_{L_2}^\dagger).
\]
Transforming the states while silently keeping the old numerical
joining map need not preserve the overlap. Independent rephasings
of the three states give
$A_{abc}\mapsto
e^{i\theta_a+i\theta_b-i\theta_c}A_{abc}$ because the parent is a
bra. Coherent phase conventions, or the invariant combinations
already constructed, are therefore needed for complex comparisons.

In the scaling picture, Euclidean time extends along three cylinders
meeting at a pair of pants. Its conformal map to a punctured sphere
relates ratios of overlaps to three-point functions. This geometric
statement presumes that the chosen microscopic seams approach the
specified continuum sewing.

\subsection{The primary geometry factor}

Put $x=L_1/L$ and define the positive entropy function
\[
 H(x)=-x\log x-(1-x)\log(1-x).
\]
An explicit cylinder coordinate is
\[
 z(w)=\frac L{2\pi}
 \bigl[x\log(w-1)+(1-x)\log w+H(x)\bigr].
\]
The incoming punctures are at $w=1,0$ and the outgoing one
is at infinity. The derivative vanishes at $w=1-x$; the
constant $H(x)$ places this joining point at zero real cylinder
time. Logarithm branches specify the seam phases.

The local disk coordinates of unit-normalized cylinder states are
$\xi_1=e^{2\pi z/L_1}$, $\xi_2=e^{2\pi z/L_2}$, and
$\xi_3=e^{-2\pi z/L}$. Near the outgoing puncture use $u=1/w$.
The minus sign is essential: the incoming ends have
$\Re z\to-\infty$, while the outgoing end has
$\Re z\to+\infty$; all three local coordinates must vanish
at their respective punctures. Normalizing states by these disk
coordinates already includes the cylinder-to-plane state factors;
one must not insert a second factor $(2\pi/L_i)^{\Delta_i}$.
Their inverse Jacobians have magnitudes
\[
 \left|\frac{dw}{d\xi_1}\right|=e^{-H/x},\qquad
 \left|\frac{dw}{d\xi_2}\right|=e^{-H/(1-x)},\qquad
 \left|\frac{du}{d\xi_3}\right|=e^H.
\]
A scalar primary of dimension $\Delta$ contributes the
$\Delta$th power of the relevant magnitude. Dividing by the
nonzero vacuum overlap with exactly the same lengths and seams
cancels the overall conformal anomaly and its gluing normalization.
Different vacuum shapes require a further shape factor, as in
the current test below.
The ideal CFT prediction for three scalar primaries in these coordinates,
to which a correctly calibrated lattice ratio should converge, is
\begin{equation}
 \frac{A_{abc}}{A_{000}}
 =e^{-H(x)[\Delta_a/x+\Delta_b/(1-x)-\Delta_c]}
       C_{ab}^{\ c}.
 \label{eq:primary-pants-factor}
\end{equation}
The coefficient is the plane correlator with the outgoing
conjugate field at infinity. For nonzero spins the same
calculation retains the holomorphic and antiholomorphic
Jacobians separately; their coordinate phases must be fixed
before comparing complex amplitudes.

At equal halves the factor becomes
$2^{-2\Delta_a-2\Delta_b+\Delta_c}$. For a normalized
conjugate pair with outgoing vacuum, $C_{a\bar a}^{\ 0}=1$,
so the ratio is
\[
 \exp\!\left[-\frac{\Delta_aH(x)}{x(1-x)}\right].
\]
This is the geometric function used in the later current test;
for spinful conjugate fields it gives the magnitude, with the
coordinate phase retained separately. At equal halves and
dimension one its value is $1/16$.
The exponential rescaling supplies an OPE estimator only if
the lattice states actually converge to those primaries.

\subsection{The vacuum size dependence}

For the normalizable identity vacuum in this periodic pants geometry,
the conformal-anomaly calculation gives
\[
 A_{000}(xL,(1-x)L)\sim C(x)L^{-c/8},\qquad 0<x<1
 \quad\text{fixed}.
\]
The microscopic seam must approach the specified sewing, with no
additional leading scaling field at the junction. The local twofold
branch has scaling dimension $\Delta_2=c/8$; Section~\ref{sec:MPS-replica}
derives that dimension from the Schwarzian. Applying it to the
normalized three-cylinder geometry is the geometric input here.
The squared overlap consequently has exponent $c/4$, and the
logarithmic bipartite fidelity is
$-\log|A_{000}|^2=(c/4)\log L-2\log|C(x)|+\cdots$.
The primary ratios above had canceled this vacuum anomaly and could
not determine it by themselves.

At two sizes with the same split and seam, the constant cancels:
\[
 c^{(\mathrm{vac})}_{L,L'}
 =-8\frac{\log|A_{000}(xL',(1-x)L')/A_{000}(xL,(1-x)L)|}
          {\log(L'/L)}.
\]
This is a velocity-free finite-size diagnostic. Corrections to the
leading geometric law remain in the ratio.

The saved equal-split vacuum overlap is $.4306101916422253$ for
$9+9\to18$, compared with the earlier $6+6\to12$ value
$.47886168385663636$. Their two-size exponent gives
$c^{(\mathrm{vac})}_{12,18}=2.0955399664$. This is an informative geometric observable;
it does not control conical, bulk, or logarithmic corrections.

\subsection{A stress overlap from the Schwarzian}

Stress is not a primary, and this gives a particularly simple pants
observable. With the same $x$ and $H(x)$, put
\[
 B(x)=x(1-x)e^{2H(x)}.
\]
For one normalized chiral stress state, the ideal pants map gives
\begin{equation}
 \left|\frac{A_{00T}}{A_{000}}\right|
       =B(x)\sqrt{c/8}.
 \label{eq:pants-stress}
\end{equation}
For example, an outgoing local coordinate can be written
$\xi=e^{-H(x)}u(1-u)^{-x}$. Reversing its power series gives
\[
 u=e^H\xi-xe^{2H}\xi^2+\frac{3x^2-x}{2}e^{3H}\xi^3
       +O(\xi^4).
\]
Thus $u'(0)=e^H$, $u''(0)=-2xe^{2H}$, and
$u'''(0)=3(3x^2-x)e^{3H}$. The Schwarzian
$u'''/u'-\tfrac32(u''/u')^2$ is
$\{u,\xi\}_{\xi=0}=-3x(1-x)e^{2H(x)}$. Inserting the stress
transformation law $T(\xi)=(du/d\xi)^2T(u)+(c/12)\{u,\xi\}$ and
dividing by $\sqrt{c/2}$ gives the signed ratio
$-B(x)\sqrt{c/8}$ in this coordinate convention, hence
\eqref{eq:pants-stress}. A rephased stress state can change the
sign; the modulus is invariant. At equal halves $B=1$.

The value comes from the Schwarzian vacuum term, not from a primary
coefficient $C(1,1,T)$. This explains why it survives even though an
ordinary plane one-point function of stress vanishes. Agreement with
the energy-Fourier estimator would be a useful common-stress test;
it is not guaranteed at accessible sizes.

\section{Connected overlap invariants}

All eight overlaps with either vacuum or one selected field on each
leg contain more information than a single three-field amplitude.
Normalize by $A_{000}$, call the resulting quantities $R$, and set
\begin{align*}
 m_1&=R_{100},\quad m_2=R_{010},\quad m_3=R_{001},\\
 B_{12}&=R_{110}-m_1m_2,
 \quad B_{13}=R_{101}-m_1m_3,
 \quad B_{23}=R_{011}-m_2m_3,\\
 K&=R_{111}-m_1R_{011}-m_2R_{101}-m_3R_{110}
       +2m_1m_2m_3.
\end{align*}
These are the algebraic analogues of connected moments. Under an
independent affine change of external field on each leg,
$X_i\mapsto u_iX_i+v_i1$, direct expansion gives
\[
 B_{ij}\mapsto u_iu_jB_{ij},\qquad
 K\mapsto u_1u_2u_3K.
\]
Thus $K^2/(B_{12}B_{13}B_{23})$ is unchanged.

For identical normalized Hermitian scalar primaries in an ideal pants geometry,
the external coordinate factors cancel between the triple coefficient
and the square root of the three pair coefficients. With consistent
real signs,
\[
 \frac{K}{\sqrt{B_{12}B_{13}B_{23}}}=C_{XXX}.
\]
The pair factors must be nonzero. For a charged order-three field,
$C_{XXX}$ can be nonzero while $\langle XX\rangle=0$; the
identical-label quotient is then undefined. One must use an
appropriate conjugate-label construction instead.
For stress, its two- and three-point normalizations instead give
\begin{equation}
 c=\frac{8B_{12}B_{13}B_{23}}{K^2}.
 \label{eq:connected-stress-c}
\end{equation}
The complex products, rather than products of absolute values, are
required. A wrong sign can turn an inconsistent finite tensor into an
apparently positive central charge.

Affine cancellation has a definite scope. It removes identity mixing,
leg rescalings, and the affine terms generated by stress coordinate
changes. A physical nonidentity insertion at the branch point produces
a four-point correlator with a connected part that need not be fixed
by the lower moments. The scalar invariant therefore does not cancel
every conical correction.

\subsection{A stronger cancellation for three stresses}

Stress has additional Ward identities that a general scalar lacks.
Assume one common local branch expansion is used in all eight
overlaps, with exact continuum vacuum or holomorphic stress on
each external leg. A cover field of dimension $w$ contributes with
relative size $L^{-w/2}$ at a twofold branch. Correlators of any
nonvacuum Virasoro primary with only external stresses vanish:
Ward recursion reduces them to the primary's zero one-point
function. The same holds for its descendants. Only the vacuum
branch family can contribute.

An extra branch stress $T(b)$ acts on a plane stress correlator by
\[
 \sum_i\left[
 \frac1{b-z_i}\partial_{z_i}
 +\frac2{(b-z_i)^2}
 +\frac{c}{2(b-z_i)^4}\mathcal R_i\right],
\]
where $\mathcal R_i$ removes the $i$th stress from that moment.
These are respectively an infinitesimal displacement of a
puncture, a leg rescaling, and an affine leg shift. The connected
invariant \eqref{eq:connected-stress-c} is invariant under all
three. Its linear response to the branch stress therefore vanishes;
differentiating in $b$ gives the same result for every
$\partial^nT(b)$.

The first new vacuum quasiprimary is
$\Lambda=:TT:-(3/10)\partial^2T$, with squared norm
$N_\Lambda=c(5c+22)/10$. Choose plane punctures $0,1,\infty$,
where the unperturbed connected pairs are $c/2$ and the triple is
$c$. Ward identities give, for a branch $\Lambda(b)$,
\begin{align*}
 \delta B_{01}&=\frac{N_\Lambda}{b^4(b-1)^4},&
 \delta B_{0\infty}&=\frac{N_\Lambda}{b^4},&
 \delta B_{1\infty}&=\frac{N_\Lambda}{(b-1)^4},\\
 \delta K&=\frac{2N_\Lambda(b^2-b+1)^2}{b^4(b-1)^4}.
\end{align*}
In $\delta\log c=\sum\delta B_{ij}/B_{ij}-2\delta K/K$,
these cancel by the elementary identity
\[
 1+(b-1)^4+b^4=2(b^2-b+1)^2.
\]
Thus even the linear weight-four correction, of nominal order
$L^{-2}$, disappears. Differentiating also removes
$\partial\Lambda$, of nominal order $L^{-5/2}$.

Because the invariant is nonlinear, this does not yet control
quadratic terms. For the deliberately incomplete branch state
$1+tT(b)$, put $q=b^{-2}$, $r=(b-1)^{-2}$. Direct substitution
of the stress Ward moments shows that the quadratic term in
$\log c$ is
\[
 (c+2)\left[-\tfrac12(q^2+r^2+q^2r^2)+qr(q+r+1)\right]=0.
\]
The same polynomial identity proves the equality. Differentiating
in the common branch location removes the mixed nonlinear
$T,\partial T$ term. Vacuum fields through weight five are
spanned by the derivatives just considered and $\Lambda$;
the only nonlinear combinations below total weight six are
$T^2$ and $T\partial T$. Hence every local branch correction
through order $L^{-5/2}$ cancels under the stated hypotheses.

Order $L^{-3}$ can survive. For example, the same bare branch
at $b=1/2$, with $c=2$, gives
$c_{\mathrm{connected}}(t)=2+6144t^3+O(t^4)$.
This is a calculation of the allowed local branch contribution,
not an error estimate for lattice data. Bulk irrelevant fields,
external-state contamination, and independently regulated branch
points are not removed by it. In particular it does not justify
fitting every finite three-stress estimator to $c+aL^{-3}$.

\begin{numericalrecord}[A useful disagreement between observables]
The scalar connected overlap invariant is $1.12049321$ at $6+6$ and
$1.12397612$ at $9+9$. Meanwhile the $L=18$ fixed range-five local
epsilon sources give absolute self cubics approximately
$.93242$--$.93755$. Both finite diagnostics are fairly stable, but they
disagree substantially if both are read as the same continuum OPE
coefficient. This is evidence that observable-dependent corrections
remain, not an invitation to average the two numbers.
\end{numericalrecord}

The finite stress tensor supplies an even sharper consistency check.
For equal halves, define the raw pair ratios
\[
 K_{13}=\frac{A_{T0T}A_{000}}{A_{T00}A_{00T}},\qquad
 K_{12}=\frac{A_{TT0}A_{000}}{A_{T00}A_{0T0}}.
\]
The ideal Ward values are $K_{13}=1-16/c$ and
$K_{12}=1+2/c$. A saved $6+6$ input-output ratio gives the
encouraging effective $c=2.07127$. But its input-input ratio is
$K_{12}=-4.21244$, which has the wrong sign relative to the
positive-charge Ward prediction. The all-eight connected formula
returns $-6.43646$ there and $-43.18719$ at $6+9$. These are
failures of the selected finite tensor to obey the continuum Ward
relations, not negative physical central charges. A separate Ising
control gives $.49938$ and $.49955$ at nearby sizes, approaching
its known $c=1/2$.

Even a correction hypothesis can be overdetermined. In the
equal-split model with only a branch insertion $1+tT(1/2)$, retain
the local-coordinate Schwarzian shifts:
\[
 X_0=T(0)+c/2,\qquad X_1=T(1)+c/2,\qquad
 X_\infty=T(\infty)-c/16.
\]
Ward identities give
\[
 \langle X_0X_1X_\infty\rangle_t
 =\frac{c(64+30c-c^2)}{64}(1+24t).
\]
For $c>0$, $c\ne32$, the nearly zero raw three-stress overlap
would require $t=-1/24$. At that value,
$K_{12}=-3-6/c$ and $K_{13}=3/2-24/c$, hence
\[
 K_{13}=4K_{12}+27/2.
\]
The measured $K_{12}$ predicts about $-3.34974$, whereas the
measured $K_{13}$ is about $-6.7247$. One arbitrary stress-branch
coefficient cannot reconcile the three measurements. Higher branch
fields, bulk corrections, and external-state errors remain separate
possibilities. The lesson is to test a correction model against
more observables than it has adjustable parameters.

\section{A current test independent of the mixing angle}

A state with spin one need not be a holomorphic current. A current
has $(h,\bar h)=(1,0)$ and must satisfy both chirality and normalized
gluing tests. A conjugate holomorphic current has the same spin;
conjugating its internal charge does not reverse its chirality.
Pairing opposite momenta can therefore test the wrong overlap.

For normalized conjugate primaries of dimension $\Delta$, the
identity OPE coefficient is one. At split $x$, the ideal pants
amplitude normalized by the same-shape vacuum overlap is
\[
 G(x;\Delta)=
 \exp\!\left[-\frac{\Delta H(x)}{x(1-x)}\right].
\]
For a current and equal halves this is $1/16$. If the vacuum
denominator is instead taken at a different shape $y$, the leading
vacuum shape factor must also be included:
\[
 G(x;1)\exp[-\tfrac c2(f(x)-f(y))],
 \qquad f(x)=\frac{F(x)+F(1-x)}{12},
\]
where $F(x)=(3-3x+2x^2)\log x/(1-x)$.
Both children must grow when using this expression asymptotically.
Keeping one child small cannot be repaired merely by inserting the correct
leading shape factor.

To allow mixing, take orthonormal bases $u_i,v_j$ of two proposed
current subspaces and form
\[
 M_{ij}=\frac{\langle0,L|\bigl(|u_i,L_1\rangle\otimes|v_j,L_2\rangle\bigr)}
 {A_{000}(L,y)}.
\]
Unitary basis changes give $M\mapsto U^TMV$. Its singular values
are invariant under state phases and arbitrary unitary mixing.
If the current and its conjugate both converge inside the retained
subspaces, the splits approach fixed fractions $x,y\in(0,1)$,
and the gluing converges as assumed, then
\[
 \liminf_{L\to\infty}\|M_L\|_2
 \geq G(x;1)\exp[-\tfrac c2(f(x)-f(y))].
\]
For asymptotically equal splits and a vacuum denominator of
the same shape, this reduces to the bound $1/16$.
This is simply the variational definition of the largest singular
value: the amplitude of any normalized pair in the two spaces is
at most $\|M\|_2$. Enlarging either space can increase this norm;
it does not show that all added states are currents.

The stored $7+8\to15$ test compares with the fixed-point geometric
value $.06162$. Its largest singular values for respectively
$1,2,3,4,6$ states per child are
\[
 .01034,\quad .01977,\quad .03395,\quad .05437,\quad .05584.
\]
Removing the projected lowest-scalar level-one descendant reduces
the six-state value to about $.03454$. These finite deficits
disfavor a current hidden in the specified small subspaces, but
finite corrections and missing subspace weight are uncontrolled.
At a smaller parent length the six-state norm exceeds the current
reference. There is no monotonic no-current theorem in these data.

A norm inequality makes the missing information explicit. Suppose
normalized current approximants lie within distances
$\delta_1,\delta_2$ of normalized vectors in the selected spaces,
and their full normalized gluing differs from the ideal value by
at most $\eta$. Literal gluing against a unit parent is a
contraction, so
\[
 \|M\|_2\geq G_{\mathrm{ideal}}-\eta
 -\frac{\delta_1+\delta_2}{|A_{000}(L,y)|}.
\]
A parent-vector error adds another analogous term. A rigorous
finite exclusion needs these errors; high state fidelity on a
single short ring does not supply their asymptotic bounds.

Finally, for an actual spin-one primary, the energy-mode convention
$H_{+1}\to\omega_L(L_{-1}+\bar L_1)$ gives
\[
 \|H_{-1}|J\rangle\|^2=\omega_L^2(\Delta_J-1),\qquad
 \|H_{+1}|J\rangle\|^2=\omega_L^2(\Delta_J+1).
\]
Thus a current has vanishing lowering norm. A nonprimary has extra
positive lowering contributions, and an unfiltered lattice norm
can also contain ultraviolet weight. A projected low-energy Ward
test, enlarged until stable, is the appropriate comparison to the
continuum identity.

\section{Perturbing the Hamiltonian: a velocity-free response test}

Take $H(\theta)=H+\theta V$ and a differentiable normalized eigenbranch.
The Hellmann--Feynman identity gives
\[
 \frac{dE_a}{d\theta}=\langle a|V|a\rangle.
\]
At an exact degeneracy one first diagonalizes $V$ in that eigenspace.
For a gap $g_a=E_a-E_0$, define
$d_a=\langle a|V|a\rangle-\langle0|V|0\rangle$. Then
\begin{equation}
 \frac{d}{d\theta}\left(\frac{g_a}{g_b}\right)
 =\frac{d_ag_b-g_ad_b}{g_b^2}.
 \label{eq:HF-ratio}
\end{equation}
Adding an identity to $V$ cancels in $d_a$. Adding $bH$ changes
$d_a$ by $bg_a$ and cancels in \eqref{eq:HF-ratio}. The response of a
gap ratio is therefore insensitive to energy shifts and overall
velocity rescalings.

For $V=+\sum_jP_j^1$, the completed $L=18$ response in coordinates
$(g_\chi/g_\phi,g_\epsilon/g_\phi)$ is approximately
$(.37336,.00731)$, whereas the $15$-to-$18$ logarithmic-size secant
is $(.03634,.01916)$. They are not parallel. Thus those two finite
drifts are not already described by a pure change of this one bare
coupling with every other correction neglected. It would be stronger
than the evidence to infer two infrared running couplings: the bare
operator can mix several fields and need not coincide with the
direction controlling size drift.

Similarly, extracting the sign of a marginal beta function from a
finite scalar self-matrix element requires an identified dimension-two
primary, a normalized scalar source, and controlled contamination by
stress and other fields. Persistent finite signs can guide a theory;
they cannot supply these missing premises.

\section{Exercises}

\begin{exercise}\label{ex:overlap-phases}
Verify all three phase invariants in \eqref{eq:triangle-invariants}
and the relation for $J^3$. Show that a symmetry conjugating all three
fields in $I_\lambda$ forces that invariant to be real even when it
attaches independent unit phases to the fields.

\par\smallskip\noindent{\footnotesize\normalfont
\hyperref[sol:ex:overlap-phases]{Solution on page~\pageref*{sol:ex:overlap-phases}.}}\par
\end{exercise}
\begin{exercise}\label{ex:overlap-descendant}
Prove by induction that $\|L_{-1}^n|h\rangle\|^2=n!(2h)_n$.
Use this to derive the first three terms of
\eqref{eq:vac-descendant-series}. Explain why the threshold
$\Delta=1/2$ is also visible from the equal-time two-point function.

\par\smallskip\noindent{\footnotesize\normalfont
\hyperref[sol:ex:overlap-descendant]{Solution on page~\pageref*{sol:ex:overlap-descendant}.}}\par
\end{exercise}
\begin{exercise}\label{ex:overlap-cumulant}
Expand $X_i\mapsto u_iX_i+v_i1$ in the definitions of $B_{ij}$ and
$K$ and prove their stated transformation laws. Identify which
assumption would fail if a fourth physical operator were inserted
at the sewing branch point.

\par\smallskip\noindent{\footnotesize\normalfont
\hyperref[sol:ex:overlap-cumulant]{Solution on page~\pageref*{sol:ex:overlap-cumulant}.}}\par
\end{exercise}
\begin{exercise}\label{ex:overlap-filter}
Prove the omitted-tail inequality for $F_L(\tau)$. Construct a
three-level example showing why small residuals of two retained
eigenvectors do not provide an energy cutoff for the omitted state.

\par\smallskip\noindent{\footnotesize\normalfont
\hyperref[sol:ex:overlap-filter]{Solution on page~\pageref*{sol:ex:overlap-filter}.}}\par
\end{exercise}
\begin{exercise}\label{ex:overlap-response}
Derive \eqref{eq:HF-ratio}. For two gap ratios collected in a vector
$R$, show why changing the normalization of $\theta$ rescales
$\partial_\theta R$ without rotating it. What further hypothesis is
needed to compare this vector with $\partial_{\log L}R$?

\par\smallskip\noindent{\footnotesize\normalfont
\hyperref[sol:ex:overlap-response]{Solution on page~\pageref*{sol:ex:overlap-response}.}}\par
\end{exercise}
\begin{exercise}\label{ex:overlap-current}
Show that unitary changes of the two incoming state bases act on
the gluing matrix as $M\mapsto U^TMV$ and preserve its singular
values. Prove the displayed finite error inequality from the
contraction property of gluing. Why does the inequality become
weaker when the vacuum overlap is small?

\par\smallskip\noindent{\footnotesize\normalfont
\hyperref[sol:ex:overlap-current]{Solution on page~\pageref*{sol:ex:overlap-current}.}}\par
\end{exercise}
\begin{exercise}\label{ex:overlap-branch}
Verify $1+(b-1)^4+b^4=2(b^2-b+1)^2$ and use it to cancel
the linear $\Lambda$ correction in the connected stress invariant.
For $q=b^{-2}$, $r=(b-1)^{-2}$, show that it also implies
$q^2+r^2+q^2r^2=2qr(q+r+1)$. Explain why two distinct branch
locations are not covered by differentiating this one-location
identity.

\par\smallskip\noindent{\footnotesize\normalfont
\hyperref[sol:ex:overlap-branch]{Solution on page~\pageref*{sol:ex:overlap-branch}.}}\par
\end{exercise}
